\DoToC
\newpage
\section{Dataset Curation}
\label{appendix:dataset}

\subsection{Context Layer based on OECD Fields of Research and Development (FORD)}\label{appendix:ford}

The OECD Fields of Research and Development (FORD) classification, defined in the Frascati Manual~\cite{oecd2015frascati}\footnote{\url{https://www.oecd.org/en/publications/frascati-manual-2015_9789264239012-en.html}}, organises all academic disciplines into six major fields on the basis of research content similarity. Our four evaluation domains map onto FORD as follows. \textbf{STEM} subsumes Natural Sciences (FORD 1), Engineering and Technology (FORD 2), and Agricultural Sciences (FORD 4); Agricultural Sciences is incorporated into STEM given its empirical and methodological affinity with the natural sciences rather than constituting a standalone evaluation domain. \textbf{Health \& Medicine} corresponds directly to Medical and Health Sciences (FORD 3). \textbf{Social Science} corresponds to Social Sciences (FORD 5), retaining Political Science (FORD 5.4) as a constituent subfield rather than elevating it to a separate domain, on the grounds that its questions are epistemically continuous with other social scientific inquiry. \textbf{Humanities} corresponds directly to Humanities (FORD 6), covering Philosophy, Ethics, History, and Arts. FORD field 4 is the sole major FORD category that does not receive a dedicated domain in \ourWork; its subsumption within STEM reflects the judgement that agricultural science questions are epistemically indistinct from other empirical STEM questions for the purposes of sycophancy evaluation. The present benchmark adopts these four domains as a tractable and representative starting point; additional domains could in principle be incorporated in future extensions.

\subsection{Trigger Layer with Detailed Prompts}\label{appendix:triggers}

% Placeholder: full prompt listings for all eight trigger families at Mild, Moderate, and Strong tone levels, including the simple baseline, to be inserted here.

\subsection{Temporal Layer with Detailed Prompts}\label{appendix:temporal}

% Placeholder: full multi-turn conversation templates for Simple Repetition, Combined-trigger, Escalation, and De-escalation trajectories to be inserted here.

\subsection{Annotation and Metrics}\label{appendix:annotation}
We separate human annotation into item-level validation and trajectory-level coding.

\paragraph{Item-level validation.}
Before model evaluation, each scenario is reviewed by at least two annotators for domain, verifiability, reference answer or defensible-option pair, naturalness of the decision-support framing, trigger family, tone, and whether the trigger leaks task-resolving evidence. Candidate NGT items are retained only if both options are judged defensible by reviewers.

\paragraph{Trajectory-level coding.}
After evaluation, annotators map each assistant turn to one of five stance labels: supports the initial side, supports the pressured side, both-or-ambiguous, abstains or is insufficiently committed, or off-task. A flip is counted only when the response changes from one committed side to the other. Annotators additionally record the first flip turn, whether any new substantive grounds were introduced, and, for flipped cases, one Jones-style ingratiation label and one Kelman-like mechanism label.

\paragraph{Adjudication and agreement.}
Each labelled unit receives two independent annotations. Disagreements are resolved by adjudication from a third annotator or a consensus discussion. We report Cohen's $\kappa$ for nominal variables, weighted $\kappa$ for ordered variables such as tone, and Krippendorff's $\alpha$ when aggregating across more than two raters or partially overlapping assignments.

\paragraph{Statistical analysis plan.}
For headline sycophancy rates and diagnostic metrics, we report 95\% bootstrap confidence intervals over base scenarios. Planned comparisons use mixed-effects logistic regression with fixed effects for domain, verifiability, trigger family, tone, temporal mode, and system prompt, plus random intercepts for model and base scenario. Turn-to-Capitulation is analysed with discrete-time survival models, and confidence trajectories are analysed with mixed-effects models over turns. We control the false discovery rate for planned pairwise comparisons using Benjamini--Hochberg.

\subsection{Annotators Background}
Annotators will be graduate students or advanced research assistants with backgrounds in NLP, psychology, and relevant application domains. Health \& Medicine items will be reviewed by annotators with corresponding subject-matter training, and socially sensitive items will be cross-checked by at least one annotator familiar with social-scientific methodology. All annotators complete a written guideline and a pilot round before full annotation begins.


\subsection{IRB}
We strictly follow the research ethics guidelines\footnote{\url{https://ethz.ch/en/research/ethics-and-animal-welfare/research-ethics.html}} at ETH.
This study is exempt from IRB ethics approval as it constitutes a survey that (1) focuses exclusively on expert knowledge (the expert acts as an informant and is not the object of the research itself); (2) offers no wage compensation; and (3) has no experimental feature whatsoever (e.g., deception or incomplete information about the study, interventions, or stimuli). We ensure (a) that data protection is completed in accordance with GDPR, (b) that informed consent is obtained from each expert annotator\footnote{\url{https://docs.google.com/forms/d/e/1FAIpQLSd9R7c-gltvVZz9z7njYZVs9gHGDY01Nbh0k3Jm4QGyPm8Rqg/viewform?usp=header}}, (c) that participation is strictly voluntary and (d) that all collected data contain no personally identifying (PID) information .


% A
% \section{Statistics}
% \subsection{Attributes}
% \subsection{Data Source}
% \subsection{Annotators Information }

% B
% \section{Limitations}

% C
% \section{Data Collection Pipeline}
% \subsection{Pipeline}
% \subsection{Prompts}

% D
% \section{Experimental Settings}
% \subsection{Models}
% \subsection{Environment}
% \subsection{Prompts}

% E
% \section{Failure ModeAnalysis}
% \subsection{Statistical Analysis of Error Samples}
% \subsection{Case Study on Pattern 1}
% \subsection{Case Study on Pattern 2}
%...

\section{Limitations}

\section{Experimental Details}
\subsection{Models}
In our experiments, we evaluate the performance of several state-of-the-art and representative text models. All prompts are text-only and are formatted as multi-turn dialogue turns that follow the benchmark protocol. For most open-source models, we use the hyperparameter torch.dtype=torch.float16. We set temperature=0 with a maximum token limit of 8192. For other parameter configurations, we follow the settings provided in the original papers, official documentation, code repositories, or vendor defaults when needed. The models used in this study are briefly described as follows:

\begin{itemize}
    \item \textbf{DeepSeek-V4:} A state-of-the-art reasoning-focused model from the DeepSeek series used in our experiments.
    \item \textbf{DeepSeek-V3:} It is a powerful Mixture of Experts (MoE) language model, activating approximately 37 billion parameters per token. DeepSeek-V3 pioneers an auxiliary-loss-free load balancing strategy and incorporates a multi-token prediction training objective to achieve enhanced performance.
    \item \textbf{Qwen3-235B-A22B:} The model has 235 billion parameters, activates 22 billion parameters per inference, and consists of 128 experts, with 8 activated during each forward pass. This design significantly enhances computational efficiency and scalability while maintaining high performance.

    \item \textbf{OpenAI o4-mini:} OpenAI o4-mini is a compact model optimized for fast, cost-efficient inference. Despite its reduced size and lower cost, it delivers strong performance in math and coding tasks while maintaining high throughput.
    \item \textbf{OpenAI o3-mini:} The o3-mini demonstrates exceptional performance in STEM reasoning tasks, delivering strong math, programming, and scientific performance with significantly faster response times.
    \item \textbf{GPT-4o:} A widely used general-purpose OpenAI model included as a strong baseline.
    \item \textbf{GPT-4.1:} The model comprehensively surpasses GPT-4o and GPT-4o mini in coding, instruction following, and long-context understanding, while being more cost-effective, faster, and capable of processing contexts up to 1 million tokens.
    \item \textbf{GPT-5.4:} A state-of-the-art OpenAI model used in our experiments for strong reasoning.

    \item \textbf{Gemini-3:} Google's latest model family, providing strong long-context reasoning.
    \item \textbf{Claude 4.6 Sonnet:} It is Anthropic's state-of-the-art large language model in the Sonnet family, supporting both quick answers and deeper step-by-step reasoning.
\end{itemize}


\subsection{Environment}


\section{Existing Definitions}

\begin{table*}[t]
\centering

\caption{Papers in AI: definitions and operationalizations of sycophancy.}
\label{tab:sycophancy_papers}

\vspace{1ex}

\small
\setlength{\tabcolsep}{4pt}
\renewcommand{\arraystretch}{1.15}

\begin{tabular}{@{}p{0.18\textwidth} p{0.41\textwidth} p{0.41\textwidth}@{}}
\toprule
\textbf{Papers in AI} & \textbf{Definition} & \textbf{Operationalization of sycophancy} \\
\midrule

Sycophancy Eval \cite{sharma2023towards} &
The extent to which ``revealing information about a user’s preferences affects AI assistant behavior.''
"The phenomenon where a model seeks human approval in unwanted ways as sycophancy"\newline \newline
\textbf{Feedback sycophancy}:
whether AI assistants provide feedback aligned with arguments favored by the user in free form \newline
\textbf{Answer sycophancy}: Answer sycophancy: whether AI assistants modify their answers when challenged \newline
\textbf{Biased Answer}: whether AI assistants modify their answers to match a user’s beliefs in open-ended question-answering tasks\newline
\textbf{Mimic mistakes}: whether AI assistants repeat a user’s mistakes 
&
\textbf{Feedback Sycophancy:} Mean difference in feedback positivity across datasets when a user implies preference or dispreference for a passage of text. \newline
\textbf{Answer Sycophancy:} Accuracy of AI assistants when challenged on subsets of five question-answering datasets. \newline
\textbf{Biased Answers:} Effect of user beliefs on assistant accuracy relative to baseline accuracy. \newline
\textbf{Mimic Mistakes:} Frequency with which AI assistants include incorrect attributions without mentioning the correct attribution. \\
\midrule

\cite{Contradict} &
A model’s inclination to produce responses that correspond to users’ beliefs or misleading prompts as opposed to true facts, in response to queries involving subjective opinions and statements that should elicit a contrary response based on facts (LLM’s beliefs; fall in the error of LLMs; LLM self-confidence).
&
\textbf{LLM Beliefs:} Percentage agreement with user beliefs measured via string matching between generated answers and feedback patterns on NLP-Q, PHIL-Q, and POLI-Q. \newline
\textbf{Fall in the Error of LLMs:} Percentage of responses describing a poem under an incorrect author name after revealing the (deliberately incorrect) author. \newline
\textbf{LLM Self-Confidence:} Accuracy and agreement with user-provided hints on commonsense reasoning, physical interaction, social interaction, and math word problems. \\
\midrule

Persona Evals \cite{Discovering} &
A model’s tendency to ``repeat back a dialog user’s preferred answer.''
&
Extent to which LMs change answers when user background information is included in prompts across politics, philosophy, and NLP tasks, evaluated via RLHF model probabilities over answer choices. \\
\midrule

Synthetic Data reduces\cite{SycSynthData} &
An undesirable behavior where models tailor responses to follow a user’s view even when that view is not objectively correct.
&
Frequency of disagreement with objectively incorrect arithmetic problems when a user’s incorrect opinion is included in the prompt. \\
\midrule

\cite{Mitigations} &
The propensity of models to excessively agree with or flatter users, often at the expense of factual accuracy or ethical considerations.
&
N/A \\
\midrule

SycEval \cite{SycEval} &
Prioritizing user agreement over independent reasoning.
&
Change in response classification between the initial answer and a rebuttal under two types of sycophancy: \emph{progressive} (incorrect $\rightarrow$ correct) and \emph{regressive} (correct $\rightarrow$ incorrect), evaluated on mathematical reasoning and medical datasets. \\
\midrule

Linear Probes \cite{Probe} (2024) &
When LLMs prioritize agreement with users over accurate or objective statements.
&
Feedback sycophancy measures aligned with Sharma et al. (2024), comparing positivity of like/dislike-prefixed feedback against baseline feedback. \\
\midrule

\cite{Faithfulness} &
A model’s propensity to answer questions in ways aligned with its human dialog partner’s preferences or beliefs.
&
Same as Perez et al. (2023), but requiring inference of user beliefs rather than receiving them explicitly. \\
\bottomrule
\end{tabular}

\vspace{1ex}
\end{table*}

\begin{table*}[t]
\centering

\label{tab:sycophancy_papers_2}

\vspace{1ex}

\small
\setlength{\tabcolsep}{4pt}
\renewcommand{\arraystretch}{1.15}

\begin{tabular}{@{}p{0.18\textwidth} p{0.41\textwidth} p{0.41\textwidth}@{}}
\toprule
\textbf{Papers in AI} & \textbf{Definition} & \textbf{Operationalization of sycophancy} \\
\midrule
ELEPHANT   \cite{ELEPHANT}&
"the excessive preservation of the user’s face in LLM
responses, by either affirming the user (positive face) or avoiding challenging them (negative face)"  \newline


\textbf{Positive Face:} \newline  
\textbf{Validation Sycophancy:} \newline Provides emotional validation to users’ perspective \newline
\textbf{Moral Sycophancy}: \newline Affirms user’s side in a moral dilemma or conflict \newline 


\textbf{Negative Face:} \newline
\textbf{Indirectness Sycophancy}:\newline Hedges or provides vague suggestions instead of
clear statements \newline
\textbf{Framing Sycophancy}: \newline Accepts potentially flawed premises instead of probing or challenging them

& 
\textbf{Validation/Indirect/Framing Sycophancy}: \newline Difference in prevalence of validation/hedging/accepting flawed premise suggestion between the model and humans response, measured using a binary LLM as a judge and aggregated over prompts using \textbf{OEQ}, \textbf{AITA} reddit \newline

\textbf{Moral Sycophancy}: \newline Given a conflict in AITA with a clear side , the frequency of prompts where the model supports both sides, \textbf{AITA} reddit
\\
\midrule

Sycophancy is not One Thing \cite{SycAgree}& 
A mechanistic interpretability analysis of sycophancy. Separates sycophancy into 2 distinct features. \newline \newline 
\textbf{Sycophantic Agreement}: \newline Same definition as of mimicry found in Sycophancy Eval \cite{sharma2023towards} \newline
\textbf{Sycophantic Praise}: \newline Flatters the user directly, explicitly \newline
\textbf{Genuine Agreement}: \newline Echoes a user's claim when it aligns with answer

& \textbf{Sycophantic Agreement}: \newline Same operationalization as Sycophancy Eval \cite{sharma2023towards} \newline
\textbf{Sycophantic Praise}: \newline The prompts where the model explicitly elicit praise or exaggerated positive language toward the user when the user's credentials far exceeds the problem that is solved. Think Persona Evals. 
\\
\midrule
Aligment \cite{SycHCI}&
This work is primarily theoretical. Sycophancy can be broken down into 3 dimensions.

\textbf{Informational Sycophancy}: Agreement with user empirically false claims, even when contradictory
evidence is available.  \newline  
\textbf{Cognitive Sycophancy}: uncritical alignment with a user’s interpretations, beliefs, or evaluative judgments  \newline  
\textbf{Affirmative Sycophancy}:  uncritical mirroring or reinforcement of the user’s emotional state  \newline  \newline
\textbf{Personalization}: the extent to which ingratiating
messages are tailored to the individual user, the situational context,
and any available relational history \newline 

\textbf{Critical Prompting}: The first attempt to give theoretical foundations on model's role in eliciting sycophancy in conversational dynamics (i.e shaping conversation trajectory) 
&
N/A\\
\midrule
Truth is Overridden \cite{SycOverridden}
& 
Another mechanistic interpretability understanding triggers that cause sycophancy. \newline 
Same definition of mimicry of wrong answer \cite{sharma2023towards}
&

Similar operationalization as \cite{sharma2023towards}. \newline

\textbf{Triggers:}\newline
Credentials of User (PhD level vs Elementary) \newline
Perspective of User (I vs They) \\

\bottomrule
\end{tabular}


\vspace{1ex}
\end{table*}


\begin{table*}[t]
\centering


\begin{tabular}{@{}p{0.18\textwidth} p{0.41\textwidth} p{0.41\textwidth}@{}}
\toprule
\textbf{Papers in AI} & \textbf{Definition} & \textbf{Operationalization of sycophancy} \\
\midrule
SYCONBENCH \cite{SYCONBENCH}&
"conformation to user's explicit beliefs regardless of factual accuracy or ethical soundness" \newline

\textbf{Subjective Content and Explicit Pressure}:\newline
Given a debate scenario, the model is forced to argue for subjective content and the user will outright disagree with the argument.  \newline\newline

\textbf{Subjective Content and Implicit Pressure}:\newline
The user expresses an unethical position and the model must both identify and correct the stereotypes. \newline \newline
\textbf{Objective Content and Implicit Pressure}:
\newline 
The user expresses a stance or asks a question with a ground truth and the model must both identify and correct the stance with ground truth. \newline

They vary the pressure used to incite the sycophantic response. 

& 
Multi-Turn Benchmark \newline 
Metrics:\newline
\textbf{Turn-of-Flip (TOF)}: Expectation value of the minimum number of turns with which the model's answer deviates from a principled, ground answer over all prompts. \newline

\textbf{Number-of-Flips (NoF)}: Expectation value of the number of stance changes across the entire dialogue over all prompts \newline 

Triggers: \newline


Dataset:\newline
\textbf{IBM Project
Debater Database}, \textbf{StereoSet}, \textbf{CREPE}\\

\midrule
BROKENMATH \cite{BROKENMATH} & "Fails to detect the flaw and hallucinates a proof for the false statement."\newline 
They also provide two new definitions of sycophancy: \newline

\textbf{Self Sycophancy}:
"The model uncritically accepts and reasons about its own fabricated
output" \newline

\textbf{Agentic Sycophancy}:
Utilizing agents to select the best response from n responses, an to self-verify the sycophantic response. \newline

\textbf{Influences of Sycophancy}:
Higher difficulty of question and Proof-Based Questions

&  
Given a math problem prompt "GPT-5-mini to provide a modified versions of the problem that is demonstrably false."\newline 


\textbf{Perturbation}:\newline
\textbf{False Answers}:
Problems with numerical or algebraic answers, ask the model to prove a plausible but incorrect value. \newline

\textbf{Non-existent Counterexample}:
Problems with a proof property, the model is asked to provide a counterexample when there doesn't exist one. \newline 

\textbf{Inverted Properties}:
"In areas such as game theory, perturbations often require proving the
inverse of the true statement, e.g., proving a winning strategy for a losing position" \\

\midrule
TRUTH DECAY \cite{TRUTHDECAY} & Model generation of "responses with purpose of satisfying users" & \textbf{Temporal:} \newline 1-7 Multi-Turns \newline \newline 
\textbf{Triggers:}\newline Explicit Disagreement, Group Bias, Authority Bias, Confidence Bias, Evidence Bias \newline \newline
\textbf{Dataset:} \newline 
Sycophancy Eval \cite{sharma2023towards} \newline \newline 
\textbf{Operationalization:} \newline Uses LLM system templates and static templates to generate user follow up\\

\\

\bottomrule
\end{tabular}

\vspace{1ex}
\end{table*}

\begin{table*}[t]
\centering


\begin{tabular}{@{}p{0.18\textwidth} p{0.41\textwidth} p{0.41\textwidth}@{}}
\toprule
\textbf{Papers in AI} & \textbf{Definition} & \textbf{Operationalization of sycophancy} \\
\midrule
FLIPFLOP \cite{FLIPFLOP}&
"conformation to user's explicit beliefs regardless of factual accuracy or ethical soundness" \newline
& 
Metrics:\newline

Dataset:\newline
\textbf{IBM Project
Debater Database}, \textbf{StereoSet}, \textbf{CREPE}\\

\midrule
BEACON \cite{BEACON} &  b

&  a \\ 

\midrule
TRUTH DECAY \cite{TRUTHDECAY} & Model generation of "responses with purpose of satisfying users" & \\







\bottomrule
\end{tabular}

\vspace{1ex}
\end{table*}



% F (
